\chapter{Analysis of Functions of More than One Variable}

Functions of more than one variable are common in many fields, including physics, engineering, and economics. In this chapter, we will explore the analysis of such functions.

\section{Functional Limits}

In order to introduce the concept of differentiability for functions of more than one variable, we first need to define some previous concepts such as functional limits. This concept requires however some topological concepts such as neighborhood, open set, boundary point, boundary, closure, and convergence. Even though these terms could have a much general definition in term of topological spaces, but we will stick to the Euclidean space $\mathbb{R}^n$, and we will give the definitions in this context.
\begin{definicion}[Neighborhood]
    Given a point $x \in \mathbb{R}^n$, a neighborhood of $x$ is, for each $\veps > 0$, the set
    \[
        N_\veps(x) = \{y \in \mathbb{R}^n : \|y - x\| < \veps\}.
    \]
\end{definicion}
\begin{definicion}[Open Set]
    A set $U \subseteq \mathbb{R}^n$ is called open if for every point $x \in U$, there exists an $\veps > 0$ such that $N_\veps(x) \subseteq U$.
\end{definicion}
\begin{prop}
    Given a collection of open sets $\{U_i\}_{i \in I}$:
    \begin{itemize}
        \item The union $\bigcup\limits_{i \in I} U_i$ is open.
        \item The intersection $\bigcap\limits_{i=1}^n U_i$ of a finite number of open sets is open.
    \end{itemize}
\end{prop}
\begin{definicion}[Complement]
    Given a set $A \subseteq \mathbb{R}^n$, the complement of $A$ is the set
    \[
        A^c := \bb{R}^n \setminus A = \{x \in \mathbb{R}^n : x \notin A\}.
    \]
\end{definicion}
\begin{definicion}[Boundary Point]
    Given a set $A \subseteq \mathbb{R}^n$, a point $x \in \mathbb{R}^n$ is called a boundary point of $A$ if:
    \begin{equation*}
        \forall \veps > 0, \qquad N_\veps(x) \cap A \neq \emptyset \quad \land \quad N_\veps(x) \cap A^c \neq \emptyset.
    \end{equation*}
\end{definicion}
\begin{definicion}[Boundary]
    Given a set $A \subseteq \mathbb{R}^n$, the boundary of $A$, denoted by $\partial A$, is the set of all boundary points of $A$.
\end{definicion}
\begin{definicion}[Closure]
    Given a set $A \subseteq \mathbb{R}^n$, the closure of $A$, denoted by $\overline{A}$, is the set
    \[
        \overline{A} := A \cup \partial A.
    \]
\end{definicion}
\begin{definicion}[Convergence]
    A sequence $\{x_n\}_{n\in \mathbb{N}}$ in $\mathbb{R}^n$ is said to converge to a point $x \in \mathbb{R}^n$ if:
    \[
        \forall \veps > 0, \qquad \exists M \in \mathbb{N} : n > M \implies x_n \in N_\veps(x).
    \]
    In this case, we write $\lim\limits_{n \to \infty} x_n = x$ or $\{ x_n \} \to x$ as $n \to \infty$.
\end{definicion}
\begin{prop}
    A sequence $\{x_n\}_{n\in \mathbb{N}}$ in $\mathbb{R}^n$ converges to a point $x \in \mathbb{R}^n$ each component converges to the corresponding component of $x$. In other words, these two statements are equivalent:
    \begin{itemize}
        \item $\{x_n\} \to x$ as $n \to \infty$.
        \item $\{x_{i}^n\} \to x_i$ as $n \to \infty$ for each $i = 1, 2, \dots, n$.
    \end{itemize}
\end{prop}

Once we have defined the concept of convergence, we can define the concept of functional limits.
\begin{definicion}[Functional Limit]
    Let $G\subset \mathbb{R}^n$ be a non-empty set, and let $f : G \to \mathbb{R}$ be a function. We say that the limit of $f$ as $x$ approaches $x_0$ exists and is equal to $L$ if, for every sequence $\{x_n\}_{n\in \mathbb{N}}$ in $G$ that converges to $x_0$, the sequence $\{f(x_n)\}_{n\in \mathbb{N}}$ converges to $L$. This is:
    \[
        \forall \{x_n\} \subseteq G, \qquad \{x_n\} \to x_0 \implies \{f(x_n)\} \to L.
    \]
    In this case, we write $\lim\limits_{x \to x_0} f(x) = L$.
\end{definicion}

\section{Continuity}

Continuity is a fundamental concept in the analysis of functions. It allows us to understand how a function behaves near a point and is essential for defining differentiability.
\begin{definicion}[Continuity]
    Let $G\subset \mathbb{R}^n$ be a non-empty set, and let $f : G \to \mathbb{R}$ be a function. We say that $f$ is continuous at a point $x_0 \in G$ if the limit of $f$ as $x$ approaches $x_0$ exists and is equal to $f(x_0)$. This is:
    \[
        \lim_{x \to x_0} f(x) = f(x_0).
    \]

    We say that $f$ is continuous on $H\subset G$ if $f$ is continuous at every point of $H$.
\end{definicion}
\begin{prop}
    Let $G\subset \mathbb{R}^n$ be a non-empty set, and let $f : G \to \mathbb{R}$ be a function. The following statements are equivalent:
    \begin{itemize}
        \item $f$ is continuous at a point $x_0 \in G$.
        \item $f_i$ is continuous at $x_0$ for each $i = 1, 2, \dots, n$.
    \end{itemize}
\end{prop}

\section{Differentiability}

Differentiability is a stronger condition than continuity. It allows us to understand how a function changes near a point and is essential for defining the concept of the derivative. Previously, we have to define some last concepts.
\begin{definicion}[Connected Set]
    A set $A \subseteq \mathbb{R}^n$ is called connected if, for any two points $x, y \in A$, there exists a polygonal path that connects $x$ and $y$ and is entirely contained in $A$.
\end{definicion}
\begin{definicion}[Domain]
    A non-empty set $G \subseteq \mathbb{R}^n$ is called a domain if it is open and connected.
\end{definicion}

\begin{teo}[Intermediate Value Theorem]
    Let $G\subset \mathbb{R}^n$ be a connected set, and let $f : G \to \mathbb{R}$ be a continuous function. Consider two points $x, y \in G$ such that $f(x) < f(y)$. Then, for every $c\in \left[f(x), f(y)\right]$, there exists a point $z \in G$ such that $f(z) = c$.
\end{teo}

\begin{definicion}[Bounded Set]
    A set $A \subseteq \mathbb{R}^n$ is called bounded if there exists a positive real number $M$ such that $\|x\| < M$ for all $x \in A$.
\end{definicion}

\begin{teo}[Extreme Value Theorem]
    Let $G\subset \mathbb{R}^n$ be a closed and bounded set, and let $f : G \to \mathbb{R}$ be a continuous function. Then, there exist points $x_{\min}, x_{\max} \in G$ such that
    \[
        f(x_{\min}) \leq f(x) \leq f(x_{\max}) \quad \text{for all } x \in G.
    \]
\end{teo}

Once those concepts are defined, we can define the concept of differentiability.
\begin{definicion}[Differentiability]
    Let $G\subset \mathbb{R}^n$ be a domain, and let $f : G \to \mathbb{R}^m$ be a function. We say that $f$ is differentiable at a point $x_0 \in G$ if there exists a linear transformation $L : \mathbb{R}^n \to \mathbb{R}^m$ such that
    \[
        \lim_{x \to x_0} \frac{\|f(x) - f(x_0) - L(x - x_0)\|}{\|x - x_0\|} = 0.
    \]
    In this case, we write $Df(x_0) = L$ and call $L$ the derivative of $f$ at $x_0$.

    We say that $f$ is differentiable on $H\subset G$ if $f$ is differentiable at every point of $H$. We then call $Df : H \to L(\mathbb{R}^n, \mathbb{R}^m)$ the derivative of $f$ on $H$.
\end{definicion}

Some properties of differentiable functions are the following:
\begin{prop}
    Let $G\subset \mathbb{R}^n$ be a domain, and let $f : G \to \mathbb{R}^m$ be a function. If $f$ is differentiable at a point $x_0 \in G$, then $f$ is continuous at $x_0$.
\end{prop}
\begin{prop}
    Let $G\subset \mathbb{R}^n$ be a domain, and let $f : G \to \mathbb{R}^m$ be a function. If $f$ is differentiable at a point $x_0 \in G$, then the derivative $Df(x_0)$ is unique.
\end{prop}
\begin{prop}
    Let $G\subset \mathbb{R}^n$ be a domain, and let $f : G \to \mathbb{R}^m$ be a function. For every $\alpha \in \mathbb{R}$, if $f$ is differentiable at a point $x_0 \in G$, then $\alpha f$ is differentiable at $x_0$ and:
    \[
        D(\alpha f)(x_0) = \alpha Df(x_0).
    \]
\end{prop}
\begin{prop}
    Let $G\subset \mathbb{R}^n$ be a domain, and let $f, g : G \to \mathbb{R}^m$ be functions. If $f$ and $g$ are differentiable at a point $x_0 \in G$, then $f + g$ is differentiable at $x_0$ and:
    \[
        D(f + g)(x_0) = Df(x_0) + Dg(x_0).
    \]
\end{prop}
\begin{prop}[Chain Rule]
    Let $G\subset \mathbb{R}^n$ be a domain, and let $f : G \to \mathbb{R}^m$. Let $g:f(G) \to \mathbb{R}^r$ be a function. If $f$ is differentiable at a point $x_0 \in G$, and $g$ is differentiable at the point $f(x_0)$, then the composition $g \circ f : G \to \mathbb{R}^r$ is differentiable at $x_0$, and:
    \[
        D(g \circ f)(x_0) = Dg(f(x_0)) \circ Df(x_0).
    \]
\end{prop}
\begin{prop}
    Let $G\subset \mathbb{R}^n$ be a domain, and let $f, g : G \to \mathbb{R}^m$ be functions. If $f$ and $g$ are differentiable at a point $x_0 \in G$, then the dot product $f \cdot g$ is differentiable at $x_0$ and:
    \[
        D(f \cdot g)(x_0) = g(x_0) \cdot Df(x_0) + f(x_0) \cdot Dg(x_0).
    \]
\end{prop}
\begin{prop}
    Let $G\subset \mathbb{R}^n$ be a domain, and let $f : G \to \mathbb{R}^m$ and $g : G \to \mathbb{R}$ be functions. If $f$ and $g$ are differentiable at a point $x_0 \in G$ such that $g(x_0)\neq 0$, then the division $\nicefrac{f}{g}$ is differentiable at $x_0$ and:
    \[
        D\left(\frac{f}{g}\right)(x_0) = \frac{g(x_0) \cdot Df(x_0) - f(x_0) \cdot Dg(x_0)}{g(x_0)^2}.
    \]
\end{prop}


\subsection{Directional Derivatives}

Let $G\subset \mathbb{R}^n$ be a domain, and let $f : G \to \mathbb{R}^m$ be a function. Given a point $x_0 \in G$ and a vector $v \in \mathbb{R}^n$ such that $\|v\| = 1$, we say that $f$ is directionally differentiable at $x_0$ in the direction of $v$ if exists a linear transformation $L : \mathbb{R}^n \to \mathbb{R}^m$ such that:
\[
    \lim_{t\to 0} \frac{\|f(x_0 + tv) - f(x_0) - L(tv)\|}{|t|} = 0.
\]
The linear transformation $L$ is called the directional derivative of $f$ at $x_0$ in the direction of $v$, and we denote it by:
\begin{equation*}
    \dfrac{\partial f}{\partial v}(x_0) = L.
\end{equation*}
\begin{observacion}
    It should be noted that the definition of directional derivative is similar to the definition of differentiability, but instead of considering the limit as $x$ approaches $x_0$ in any direction, we only consider when $x$ approaches $x_0$ along the direction of the vector $v$, i.e., when $x = x_0 + tv$ for some scalar $t$ that approaches $0$. Therefore, if a function is differentiable at a point, then it is also directionally differentiable at that point in every direction.
\end{observacion}

In the practice, the directional derivative is just a normal derivative of a single variable function. In fact, if we define the function:
\Func{\varphi_v}{\left]-\delta, \delta\right[}{\mathbb{R}^m}{t}{f(x_0 + tv)}
where $\delta > 0$ is such that $x_0 + tv \in G$ for all $t \in \left]-\delta, \delta\right[$, then the directional derivative of $f$ at $x_0$ in the direction of $v$ is just the derivative of $\varphi_v$ at $t = 0$, i.e.:
\begin{equation*}
    \dfrac{\partial f}{\partial v}(x_0) = \varphi_v'(0).
\end{equation*}
where, if $m>1$, the derivative $\varphi_v'(0)$ is just the vector of the derivatives of each component of $\varphi_v$ at $t = 0$.

\subsection{Partial Derivatives}

Let $G\subset \mathbb{R}^n$ be a domain, and let $f : G \to \mathbb{R}^m$ be a function. Given a point $x_0 \in G$ and an index $i \in \{1, 2, \dots, n\}$, we say that $f$ is partially differentiable at $x_0$ with respect to the variable $x_i$ if it is directionally differentiable at $x_0$ in the direction of the standard basis vector $e_i$, where $e_i$ is the vector with a $1$ in the $i$-th position and $0$ in all other positions. The partial derivative of $f$ at $x_0$ with respect to the variable $x_i$ is denoted by:
\begin{equation*}
    \frac{\partial f}{\partial x_i}(x_0) = \dfrac{\partial f}{\partial e_i}(x_0).
\end{equation*}

\subsection{Jacobian Matrix}

Let $G\subset \mathbb{R}^n$ be a domain, and let $f : G \to \mathbb{R}^m$ be a function. If $f$ is differentiable at a point $x_0 \in G$, then the derivative $Df(x_0)$, as a linear transformation from $\mathbb{R}^n$ to $\mathbb{R}^m$, can be represented by a matrix. This matrix is called the Jacobian matrix of $f$ at $x_0$, and it is denoted by $Jf(x_0)\in \cc{M}_{m \times n}(\mathbb{R})$.
\begin{prop}
    Let $G\subset \mathbb{R}^n$ be a domain, and let $f : G \to \mathbb{R}^m$ be a function. If $f$ is differentiable at a point $x_0 \in G$, then the Jacobian matrix of $f$ at $x_0$ is given by:
    \[
        Jf(x_0) = \begin{pmatrix}
            \dfrac{\partial f_1}{\partial x_1}(x_0) & \dfrac{\partial f_1}{\partial x_2}(x_0) & \cdots & \dfrac{\partial f_1}{\partial x_n}(x_0) \\
            \dfrac{\partial f_2}{\partial x_1}(x_0) & \dfrac{\partial f_2}{\partial x_2}(x_0) & \cdots & \dfrac{\partial f_2}{\partial x_n}(x_0) \\
            \vdots & \vdots & \ddots & \vdots \\
            \dfrac{\partial f_m}{\partial x_1}(x_0) & \dfrac{\partial f_m}{\partial x_2}(x_0) & \cdots & \dfrac{\partial f_m}{\partial x_n}(x_0)
        \end{pmatrix}
        = \left( \dfrac{\partial f_i}{\partial x_j}(x_0) \right)_{1 \leq i \leq m, 1 \leq j \leq n}.
    \]
\end{prop}

If $n=1$, i.e., if $G\subset \mathbb{R}$, then the Jacobian matrix of $f$ at $x_0$ is just a column vector, and it is given by:
\[
    Jf(x_0) = \begin{pmatrix}
        f_1'(x_0) \\
        f_2'(x_0) \\
        \vdots \\
        f_m'(x_0)
    \end{pmatrix}
\]
where $f_i'(x_0)$ is the derivative of the $i$-th component of $f$ at $x_0$. We call this the derivative vector of $f$ at $x_0$.\\

If $m=1$, i.e., if $f : G \to \mathbb{R}$ is a scalar-valued function, then the Jacobian matrix of $f$ at $x_0$ is just a row vector, and it is given by:
\[
    Jf(x_0) = \begin{pmatrix}
        \dfrac{\partial f}{\partial x_1}(x_0) & \dfrac{\partial f}{\partial x_2}(x_0) & \cdots & \dfrac{\partial f}{\partial x_n}(x_0)
    \end{pmatrix}
\]
This is the gradient of $f$ at $x_0$, and we denote it by $\nabla f(x_0)=\text{grad } f(x_0)=Jf(x_0)$. It should be noted that the gradient of a scalar-valued function is a row vector that points in the direction of the greatest rate of increase of the function, and its magnitude represents the rate of increase in that direction.



\subsection{Total Differentiability}

The total differentiability of a function lets us understand how a does the value of a function change when we change all the variables at the same time. The term $dx$ represents a small change in the input variables, and the term $df$ represents the corresponding change in the output of the function. Considering $G$ a domain, and $f : G \to \mathbb{R}^m$ a function, its total differential at a point $x_0 \in G$ is given by:
\[
    df = Jf(x_0) \cdot dx
\]f scalar-valued functions, i.e., when $m=1$, where the total differential is given by:
\[
    df = \nabla f(x_0) \cdot dx = \sum_{i=1}^n \frac{\partial f}{\partial x_i}(x_0) dx_i
\]


% // TODO: Añadir red neuronal

\section{Integration of Functions of More than One Variable}

Integration of functions of more than one variable is a generalization of the concept of integration for functions of a single variable. It allows us to compute the volume under a surface, the mass of an object, and many other quantities that depend on multiple variables.

\subsection{Curve Integrals}

Let $[a, b]$ be a closed interval in $\mathbb{R}$, and let $\varphi : [a, b] \to \mathbb{R}^n$ be a continuous function. We say that $\varphi$ is the parametrization of a curve $C\subset \mathbb{R}^n$ given by:
\[
    C = \{\varphi(t) : t \in [a, b]\}.
\]
\begin{observacion}
    It should be noted that the curve $C$ is the image of the parametrization $\varphi$, and it is a subset of $\mathbb{R}^n$. A curve may have multiple parametrizations. For example, the curve $C$ given by the unit circle in $\mathbb{R}^2$ can be parametrized by $\varphi(t) = (\cos t, \sin t)$ for $t \in [0, 2\pi]$, or by $\psi(t) = (\cos(2t), \sin(2t))$ for $t \in [0, \pi]$. Both $\varphi$ and $\psi$ represent the same curve $C$, which is the unit circle.
\end{observacion}

The parameterization $\varphi$ is said to be \emph{regular} if $\varphi$ is differentiable and $\varphi'(t) \neq 0$ for all $t \in \left]a, b\right[$. In this case, we say that $C$ is a regular curve.
\begin{observacion}
    A parameterization $\varphi$ can be seen as how we traverse the curve $C$. The parameter $t$ can be thought of as time, and $\varphi(t)$ gives us the position on the curve at time $t$. The vector $\varphi'(t)$ represents the velocity of the traversal at time $t$. If $\varphi'(t) \neq 0$ for all $t$, it means that we are always moving along the curve and never stopping, which is why we call it a regular parameterization. If $\varphi'(t) = 0$ for some $t$, it means that we are momentarily stopping at that point on the curve, which can lead to complications (as the direction could change abruptly), and that's why we don't consider it regular.
\end{observacion}

\subsubsection{Length of a Curve}

The length of a curve $C$ is given by the supremum of the lengths of polygonal paths inscribed in $C$. If that supremum is finite, we say that $C$ is rectifiable, and we denote its length by $\ell(C)$.
\begin{prop}
    Let $[a, b]$ be a closed interval in $\mathbb{R}$, and let $\varphi : [a, b] \to \mathbb{R}^n$ be a regular parameterization of a curve $C$. Then, $C$ is rectifiable. If $\varphi$ is injective, then:
    \[
        \ell(C) = \int_a^b \|\varphi'(t)\| dt.
    \]
\end{prop}
% // TODO: Memorizar longitud curva


\subsubsection{Curve Integral of First Kind}
% // TODO: Memorizar Integral 1er Tipo
Let $G$ be a \emph{domain} in $\mathbb{R}^n$ (i.e., an open and connected set), and let $f : G \to \mathbb{R}$ be a \emph{continuous} function. Let $C$ be a curve in $G$ with a \emph{regular} parameterization $\varphi : [a, b] \to G$. The curve integral of first kind of $f$ along $\varphi$ is defined as:
\[
    \int_\varphi f = \int_a^b f(\varphi(t)) \|\varphi'(t)\| dt.
\]
\begin{observacion}
    The curve integral of first kind can be thought of a ``weighted'' length of the curve $C$, where the weight at each point $\varphi(t)$ is given by the value of the function $f$ at that point. If $f$ is constant, then the curve integral of first kind reduces to the length of the curve multiplied by that constant.
\end{observacion}

\begin{prop}[Linearity of the Curve Integral of First Kind]
    Let $G$ be a domain in $\mathbb{R}^n$, and let $f, g : G \to \mathbb{R}$ be continuous functions. Let $C$ be a curve in $G$ with a regular parameterization $\varphi : [a, b] \to G$. Then:
    \begin{itemize}
        \item For every $\alpha \in \mathbb{R}$, we have $\displaystyle \int_\varphi \alpha f = \alpha \int_\varphi f$.
        \item $\displaystyle \int_\varphi (f + g) = \int_\varphi f + \int_\varphi g$.
    \end{itemize}
\end{prop}
\begin{prop}[Orientation of the Curve Integral of First Kind]
    Let $G$ be a domain in $\mathbb{R}^n$, and let $f : G \to \mathbb{R}$ be a continuous function. Let $C$ be a curve in $G$ with a regular parameterization $\varphi : [a, b] \to G$. If we reverse the orientation of the curve by considering the parameterization $\psi : [a, b] \to G$ given by $\psi(t) = \varphi(a + b - t)$, then:
    \[
        \int_\psi f = \int_\varphi f.
    \]
\end{prop}
\begin{prop}[Modulus of the Curve Integral of First Kind]
    Let $G$ be a domain in $\mathbb{R}^n$, and let $f : G \to \mathbb{R}$ be a continuous function. Let $C$ be a curve in $G$ with a regular parameterization $\varphi : [a, b] \to G$. Then:
    \[
        \left|\int_\varphi f\right| \leq \int_\varphi |f|.
    \]
\end{prop}

Some important applications of the curve integral of first kind are the following:
\begin{itemize}
    \item Let's suppose that we have a wire in the shape of a curve $C$ in $\mathbb{R}^3$, and we want to compute the mass of that wire. If we know the density of the wire at each point, which is given by a continuous function $f : C \to \mathbb{R}$, then the mass of the wire can be computed as the curve integral of first kind of $f$ along a regular parameterization of $C$. That is:
    \[
        \text{Mass} = \int_\varphi f.
    \]

    \item In order to calculate ``fence areas'' we can use the curve integral of first kind. Given a domain $G$ in $\mathbb{R}^2$, a curve $C$ in $G$ with a regular inyective parameterization $\varphi : [a, b] \to G$, and a continuous function $f : G \to \mathbb{R}$ such that it represents the height of a fence at each point of $G$, then the area of the fence can be computed as the curve integral of first kind of $f$ along $\varphi$. That is:
    \[
        \text{Area} = \int_\varphi f.
    \]
    Mathematically, we compute the area of:
\[
    \{(\varphi_1(t), \varphi_2(t), z) : z\in [0, f(\varphi(t))],\ t \in [a, b]\}
\]
\end{itemize}

\subsubsection{Curve Integral of Second Kind}
% // TODO: Memorizar Integral 2º Tipo
Let $G$ be a \emph{domain} in $\mathbb{R}^n$, and let $f : G \to \mathbb{R}^n$ be a \emph{continuous} function. Let $C$ be a curve in $G$ with a \emph{regular} parameterization $\varphi : [a, b] \to G$. The curve integral of second kind of $f$ along $\varphi$ is defined as:
\[
    \int_\varphi f = \int_a^b \left(f(\varphi(t)) \cdot \varphi'(t)\right) dt.
\]

Once important aspect to take into account is the tangential component of the vector field $f$ along the parameterization $\varphi$. The tangential component of $f$ along $\varphi$ at a point $\varphi(t)$ is given by the projection of $f(\varphi(t))$ onto the direction of $\varphi'(t)$, which is given by:
\[
    \text{proj}_{\varphi'(t)} f(\varphi(t)) =
    \|f(\varphi(t))\| \cos(\alpha(t))
\]
where $\alpha(t)$ is the angle between the vectors $f(\varphi(t))$ and $\varphi'(t)$. Therefore, the curve integral of second kind can be thought of as the integral of the tangential component of $f$ along the parameterization $\varphi$.

\begin{prop}[Linearity of the Curve Integral of Second Kind]
    Let $G$ be a domain in $\mathbb{R}^n$, and let $f, g : G \to \mathbb{R}^n$ be continuous functions. Let $C$ be a curve in $G$ with a regular parameterization $\varphi : [a, b] \to G$. Then:
    \begin{itemize}
        \item For every $\alpha \in \mathbb{R}$, we have $\displaystyle \int_\varphi \alpha f = \alpha \int_\varphi f$.
        \item $\displaystyle \int_\varphi (f + g) = \int_\varphi f + \int_\varphi g$.
    \end{itemize}
\end{prop}

\begin{prop}[Orientation of the Curve Integral of Second Kind]
    Let $G$ be a domain in $\mathbb{R}^n$, and let $f : G \to \mathbb{R}^n$ be a continuous function. Let $C$ be a curve in $G$ with a regular parameterization $\varphi : [a, b] \to G$. If we reverse the orientation of the curve by considering the parameterization $\psi : [a, b] \to G$ given by $\psi(t) = \varphi(a + b - t)$, then:
    \[
        \int_\psi f = -\int_\varphi f.
    \]
\end{prop}

\begin{prop}[Modulus of the Curve Integral of Second Kind]
    Let $G$ be a domain in $\mathbb{R}^n$, and let $f : G \to \mathbb{R}^n$ be a continuous function. Let $C$ be a curve in $G$ with a regular parameterization $\varphi : [a, b] \to G$. Then:
    \[
        \left|\int_\varphi f\right| \leq \int_\varphi \|f\|.
    \]
\end{prop}


Some important applications of the curve integral of second kind are the following:
\begin{itemize}
    \item In order to calculate the area enclosed by a curve $C$ in $\mathbb{R}^2$, we can use the curve integral of second kind. Given a domain $G$ in $\mathbb{R}^2$, a curve $C$ in $G$ with a regular parameterization $\varphi : [a, b] \to G$ that traverses the curve only once in a counterclockwise direction, the area enclosed by the curve $C$ can be computed as:
    \[
        \text{Area} = \frac{1}{2} \int_\varphi (-y,x)
    \]
    \begin{itemize}
        \item If enclosing the area in a clockwise direction needs $k$ regular parameterizations $\varphi_1, \varphi_2, \dots, \varphi_k$, then the area can be computed as:
        \[
            \text{Area} = \frac{1}{2} \sum_{i=1}^k \int_{\varphi_i} (-y,x).
        \]
    \end{itemize}

    \item In order to calculate the work done by a force field $f : G \to \mathbb{R}^n$ along a curve $C$ in $G$, we can use the curve integral of second kind. Given a domain $G$ in $\mathbb{R}^n$, a curve $C$ in $G$ with a regular parameterization $\varphi : [a, b] \to G$, and a continuous function $F : G \to \mathbb{R}^n$ that represents the force field, then the work done by the force field along the curve can be computed as:
    \[
        W = \int_\varphi F.
    \]
\end{itemize}


\subsection{Potential fields}

Given a domain $G$ in $\mathbb{R}^n$, and a continuous vector field $F : G \to \mathbb{R}^n$, we say that $F$ is a potential field if there exists a differentiable scalar-valued function $\varphi : G \to \mathbb{R}$ such that:
\[
    F = \nabla \varphi = \left( \frac{\partial \varphi}{\partial x_1}, \frac{\partial \varphi}{\partial x_2}, \dots, \frac{\partial \varphi}{\partial x_n} \right).
\]
In this case, we say that $\varphi$ is a potential function for the potential field $F$. For instance, the electric field generated by a point charge is a potential field, and its potential function is the electric potential:
\[
    \vec{E} = -\nabla V
\]

Potential fields have the following important property:
\begin{prop}
    Let $G$ be a domain in $\mathbb{R}^n$, and let $F : G \to \mathbb{R}^n$ be a continuous vector field. If $F$ is a potential field with potential function $\varphi : G \to \mathbb{R}$, then for every curve $C$ in $G$ with a regular parameterization $\psi : [a, b] \to G$, we have:
    \[
        \int_\psi F = \varphi(\psi(b)) - \varphi(\psi(a)).
    \]
\end{prop}
That is, the curve integral of second kind of a potential field along a curve $C$ depends only on the endpoints of the curve, and not on the particular path taken between those endpoints. In order to find a potential function for a given vector field, we can use the following proposition.
\begin{prop}
    Let $G$ be a domain in $\mathbb{R}^n$, and let $F : G \to \mathbb{R}^n$ be a continuous vector field. If a curve integral of second kind of $F$ along any curve $C$ in $G$ depends only on the endpoints of the curve, then $F$ is a potential field. That is, for every two points $x_1, x_2 \in G$ and every regular parameterization $\psi : [a, b] \to G$ such that $\psi(a) = x_1$ and $\psi(b) = x_2$, the following integral depends only on the endpoints $x_1$ and $x_2$, i.e.:
\[
        \int_\psi F
\]
    depends only on $x_1$ and $x_2$ and not on the particular parameterization $\psi$.
\end{prop}

Even though this proposition provides us a theoretical way to prove that a vector field is in fact a potential field, in practice determining if a curve integral is path independent or not is not so easy. One condition that every potential field must satisfy is the following.
\begin{prop}
    Let $G$ be a domain in $\mathbb{R}^n$, and let $F : G \to \mathbb{R}^n$ be a continuous vector field. If $F$ is a potential field with potential function $\varphi : G \to \mathbb{R}$, then the partial derivatives of $\varphi$ satisfy the following conditions:
    \[
        \frac{\partial F_i}{\partial x_j} = \frac{\partial F_j}{\partial x_i} \quad \text{for all } i, j \in \{1, 2, \dots, n\}.
    \]
\end{prop}
This condition is necessary for a vector field to be a potential field, but it is not sufficient. In order to apply it in the opposite direction, we need to add some extra conditions on the domain $G$: it has to be simply connected, which means that it has no holes. We will however not go into the details of simply connected domains, as star-shaped domains are simply connected.
\begin{definicion}[Star-shaped]
    A set $G\subset\mathbb{R}^n$ is called star-shaped from a point $x_0 \in G$ if for every point $x \in G$, the line segment connecting $x_0$ and $x$ is entirely contained in $G$. That is, for every $x \in G$ and every $t \in [0, 1]$, we have:
    \[
        (1 - t)x_0 + tx \in G.
    \]

    We say that $G$ is a star-shaped domain if there exists a point $x_0 \in G$ such that $G$ is star-shaped from $x_0$.
\end{definicion}
\begin{definicion}[Convex]
    A set $G\subset\mathbb{R}^n$ is called convex if for every two points $x, y \in G$, the line segment connecting $x$ and $y$ is entirely contained in $G$. That is, for every $x, y \in G$ and every $t \in [0, 1]$, we have:
    \[
        (1 - t)x + ty \in G.
    \]
\end{definicion}

It is easy to see that every convex set is star-shaped, which will be useful for us, as we will see in the next proposition.
\begin{prop}
    Let $G$ be a star-shaped domain in $\mathbb{R}^n$, and let $F : G \to \mathbb{R}^n$ be a continuous vector field. If the partial derivatives of $F$ satisfy the following conditions:
    \[
        \frac{\partial F_i}{\partial x_j} = \frac{\partial F_j}{\partial x_i} \quad \text{for all } i, j \in \{1, 2, \dots, n\},
    \]
    then $F$ is a potential field.
\end{prop}

This is a very useful result, as it provides us a way to determine if a vector field is a potential field or not by just checking the equality of the mixed partial derivatives. In practice, this is much easier than trying to determine if a curve integral is path independent or not. Once we have proved that a vector field is a potential field, there are two main ways to find a potential function for that vector field:
\begin{enumerate}
    \item An origin point $x_0$ is chosen in the domain $G$ to serve as a reference point. For any point $x$ in the domain, a curve $\psi$ is defined that connects $x_0$ to $x$:
    \Func{\psi}{[0, 1]}{G}{t}{x_0 + t(x - x_0)}
    The potential function $\varphi$ is then defined as the curve integral of second kind of the vector field $F$ along the curve $\psi$:
    \[
        \varphi(x) = \int_\psi F.
    \]

    Ideally, the origin point $x_0$ could be chosen as the origin of the coordinate system, i.e., $x_0 = (0, 0, \dots, 0)$, so that the curve $\psi$ would be even simpler.

    \item The potential function $\varphi$ can be directly computed by integrating the components of the vector field $F$.
\end{enumerate}

\subsection{Multiple Integrals}

In this section, we will introduce the concept of multiple integrals, which are a generalization of the concept of integration for functions of a single variable. Multiple integrals allow us to compute volumes, masses, and other quantities that depend on multiple variables. In order to introduce them, normal areas need to be defined.
\begin{definicion}[Normal Area]
    A normal area in $\mathbb{R}^n$ is inductively defined as follows:
    \begin{itemize}
        \item A normal area in $\mathbb{R}$ is just a closed interval $[a, b]$.
        \item A normal area in $\mathbb{R}^n$ for $n > 1$ is a set of the form:
        \[
            A = \{(x_1,\dots,x_{n-1},x_n) \in \mathbb{R}^n : (x_1,\dots,x_{n-1}) \in A',\ g_1(x_1,\dots,x_{n-1}) \leq x_n \leq g_2(x_1,\dots,x_{n-1})\}
        \]
        where $A'$ is a normal area in $\mathbb{R}^{n-1}$, and $g_1, g_2 : A' \to \mathbb{R}$ are continuous functions such that:
        \[
            g_1(x_1,\dots,x_{n-1}) \leq g_2(x_1,\dots,x_{n-1}) \quad \text{for all } (x_1,\dots,x_{n-1}) \in A'.
        \]

        In that case, $A$ is called a normal area in $\mathbb{R}^n$ with respect to the variable $x_n$. The functions $g_1$ and $g_2$ are called the lower and upper bounds of the normal area $A$, respectively. The set $A'$ is called the projection of $A$ onto the first $n-1$ variables, and it is denoted by $\pi(A)$.
    \end{itemize}
\end{definicion}

Once normal areas have been defined, we can introduce the concept of a parameter-dependent integral, which is a generalization of the concept of integration for functions of a single variable. Given a normal area $A$ in $\mathbb{R}^n$ with respect to the variable $x_n$, and a continuous function $f : A \to \mathbb{R}$, the parameter-dependent integral of $f$ over $A$ is defined as:
\Func{F}{\pi(A)}{\mathbb{R}}{(x_1,\dots,x_{n-1})}{\int_{g_1(x_1,\dots,x_{n-1})}^{g_2(x_1,\dots,x_{n-1})} f(x_1,\dots,x_{n-1},x_n)\ dx_n}

In other words, the parameter-dependent integral of $f$ over $A$ is a function that assigns to each point in the projection $\pi(A)$ the value of the integral of $f$ with respect to the variable $x_n$ over the interval defined by the lower and upper bounds of $A$ at that point. In the case of a normal area in $\mathbb{R}^2$ with respect to the variable $y$, the parameter-dependent integral of $f$ over $A$ is given by:
\[
    F(x) = \int_{g_1(x)}^{g_2(x)} f(x,y)\ dy
\]

If $g_1$ and $g_2$ are differentiable and $f$ is partially differentiable with respect to $x$, then $F$ is also differentiable with:
\[
    F'(x) = \int_{g_1(x)}^{g_2(x)} \frac{\partial f}{\partial x}(x,y)\ dy + f(x, g_2(x)) \cdot g_2'(x) - f(x, g_1(x)) \cdot g_1'(x).
\]

In order to integrate a function over a normal area, we can use the following proposition.
\begin{prop}
    Let $A$ be a normal area in $\mathbb{R}^n$ with respect to the variable $x_n$, and let $f : A \to \mathbb{R}$ be a continuous function. If the parameter-dependent integral of $f$ over $A$ is integrable over the projection $\pi(A)$, then $f$ is integrable over $A$, and we have:
\[
    \int_A f = \int_{\pi(A)} F.
\]
\end{prop}
This proposition allows us to compute the integral of a function over a normal area by first computing the parameter-dependent integral of the function over the normal area, and then integrating that parameter-dependent integral over the projection of the normal area. This is a very useful result, as it allows us to compute multiple integrals by iteratively applying single-variable integration. Some important applications of multiple integrals are the following:
\begin{itemize}
    \item In 2D, the area of a region $R$ in $\mathbb{R}^2$ that is below a continuos function $f:[a, b] \to \mathbb{R}^+_0$ is the following integral:
    \[
        \text{Area} = \int_a^b f(x)\ dx.
    \]
    \item The previous well-known example can be easily extended to 3D, where the volume of a region $R$ in $\mathbb{R}^3$ that is below a continuos function $f : A \to \mathbb{R}^+_0$ defined on a normal area $A$ in $\mathbb{R}^2$ is given by:
    \[
        \text{Volume} = \int_A f.
    \]
    \item Given a set $M\subset \bb{R}^3$ with volume $V$ and homogeneous density, the center of mass of $M$, $(s_1, s_2, s_3)$, is given by:
    \begin{equation*}
        s_k = \frac{1}{V} \int_M x_k\ dx_1 dx_2 dx_3 \quad \text{for } k = 1, 2, 3.
    \end{equation*}
\end{itemize}


\subsubsection{Changing the Coordinate System}

Until know, we have just worked with the cartesian coordinate system. However, we can also work with other coordinate systems, such as the polar coordinate system in $\bb{R}^2$ or the cylindrical and spherical coordinate systems in $\bb{R}^3$. In order to change the coordinate system, some advanced mathematical tools are needed, such as the concept of ``difeomorphism''. However, the following intuitive idea can be followed:
\begin{itemize}
    \item In polar coordinates, $dA=dx\ dy$ is replaced by:
    \[
        dA = r\ dr\ d\theta
    \]
    \item In cylindrical coordinates, $dV=dx\ dy\ dz$ is replaced by:
    \[
        dV = r\ dr\ d\theta\ dz
    \]
    \item In spherical coordinates, $dV=dx\ dy\ dz$ is replaced by:
    \[
        dV = \rho^2 \sin\phi\ d\rho\ d\phi\ d\theta
    \]
    where $\phi$ is the angle between the positive $z$-axis and the line connecting the origin to the point.
\end{itemize}